% =====================================================================
%  ÔN TẬP — Phân tích dữ liệu thông minh
%  Bản rút gọn của toàn bộ thư mục on-tap/, soạn từ book của thầy:
%  https://lenbui.github.io/bookda/
%
%  Biên dịch: pdflatex on-tap.tex   (chạy 2 lần cho mục lục đúng)
%  Không cần hình, không cần file phụ. Upload mỗi file này lên Overleaf
%  là compile được.
%
%  Preamble lấy nguyên từ bao-cao/tex/main.tex vì bộ T5 + vietnamese đó
%  đã compile thành công rồi — không đổi để khỏi dò lại lỗi font.
% =====================================================================
\documentclass[11pt]{article}

\usepackage[utf8]{inputenc}
\usepackage[T5]{fontenc}
\usepackage[vietnamese]{babel}
\usepackage[a4paper,margin=2.2cm]{geometry}
\usepackage{amsmath}
\usepackage{amssymb}
\usepackage{enumitem}
\usepackage{booktabs}
\usepackage{array}
\usepackage{hyperref}

\hypersetup{unicode=true, hidelinks,
  pdftitle={On tap - Phan tich du lieu thong minh}}
\setlist{leftmargin=*, topsep=2pt, itemsep=1pt}
\emergencystretch=3em
\setlength{\parindent}{0pt}
\setlength{\parskip}{4pt}

% Nhãn "bẫy" dùng lại nhiều lần.
\newcommand{\bay}{\textbf{Bẫy.}\ }
\newcommand{\vd}{\textbf{Ví dụ.}\ }

\title{\vspace{-1.5cm}\bfseries Ôn tập — Phân tích dữ liệu thông minh}
\author{Bản rút gọn của \texttt{on-tap/}, soạn từ \url{https://lenbui.github.io/bookda/}}
\date{}

\begin{document}
\maketitle
\vspace{-1cm}

\begin{center}
\begin{tabular}{@{}ll@{}}
\toprule
\textbf{Chủ đề} & \textbf{Chương trong book} \\
\midrule
Mô tả dữ liệu, phương sai, Pearson vs Spearman & 3 \\
Khoảng tin cậy & 5 \\
Kiểm định giả thuyết, miền bác bỏ & 6 \\
Mô hình đồ thị, confounder / mediator & 7 \\
Phân tích nhân quả & 8 \\
Hồi quy tuyến tính & 9 \\
Hồi quy đa biến và logistic & 10 \\
Monte Carlo & 11 \\
\bottomrule
\end{tabular}
\end{center}

\textbf{Ba chỗ book không phủ mà vẫn nên biết.}
(1)~Chương 5 chỉ dựng khoảng tin cậy cho \emph{tỉ lệ}, không nhắc phân phối $t$,
không có khoảng tin cậy cho trung bình --- nhưng chương 6 lại dùng $t$ cho trung
bình, nên phần đó được bổ sung ở mục~\ref{sec:ktc}.
(2)~Chương 8 dừng ở ATE và RCT, không có backdoor, không có \texttt{do(\ )}, không
có điểm xu hướng; ba thứ đó chỉ cần cho vấn đáp đồ án.
(3)~Chương 7 dạy ba cấu trúc chuỗi / rẽ nhánh / hội tụ rồi để người đọc tự suy ra
tên gọi confounder và mediator, trong khi đề thi hỏi thẳng tên --- mục~\ref{sec:dothi}
gắn tên vào từng cấu trúc.

\textbf{Thứ tự ôn theo tỉ lệ điểm trên thời gian:}
kiểm định (nặng nhất, đề gần như chắc chắn rút từ bảng ở mục~\ref{sec:kiemdinh})
$\to$ hồi quy tuyến tính (bài rất mẫu mực, được điểm nhanh)
$\to$ khoảng tin cậy (dùng lại đúng bộ $z^{*}$, ôn chung cho tiết kiệm)
$\to$ mô hình đồ thị (rẻ điểm nhất, không phải tính toán)
$\to$ phần còn lại.

\tableofcontents

\newpage

%=====================================================================
\section{Mô tả dữ liệu (chương 3)}

\subsection*{Công thức}

\[
\bar{x} = \frac{1}{n}\sum_{i=1}^{n} x_i
\qquad
s^2 = \frac{\sum_{i=1}^{n}(x_i-\bar{x})^2}{n-1}
\qquad
s = \sqrt{s^2}
\]

\[
IQR = Q_3 - Q_1
\qquad
\text{râu trên} = Q_3 + 1{,}5\,IQR
\qquad
\text{râu dưới} = Q_1 - 1{,}5\,IQR
\]

Tính nhanh khi số xấu:
$\displaystyle\sum (x_i-\bar{x})^2 = \sum x_i^2 - \frac{\left(\sum x_i\right)^2}{n}$

\textbf{Ba hệ số tương quan.}

\[
r_{\text{Pearson}} = \frac{S_{xy}}{\sqrt{S_{xx}S_{yy}}}
\qquad
\rho_{\text{Spearman}} = 1 - \frac{6\sum d_i^2}{n(n^2-1)}
\qquad
\tau_{\text{Kendall}} = \frac{C-D}{n(n-1)/2}
\]

với $S_{xy}=\sum(x_i-\bar x)(y_i-\bar y)$, $S_{xx}=\sum(x_i-\bar x)^2$,
$S_{yy}=\sum(y_i-\bar y)^2$; $d_i$ là hiệu \emph{hạng}; $C$ và $D$ là số cặp đồng
hạng và nghịch hạng.

\subsection*{Chọn hệ số nào}

\begin{center}
\begin{tabular}{@{}lll@{}}
\toprule
 & \textbf{Pearson} $r$ & \textbf{Spearman} $\rho$ \\
\midrule
Đo & tuyến tính & đơn điệu (không cần thẳng) \\
Dữ liệu & định lượng & định lượng \emph{hoặc thứ bậc} \\
Ngoại lệ & rất nhạy & bền \\
$y=x^3$ & $r<1$ & $\rho=1$ chính xác \\
\bottomrule
\end{tabular}
\end{center}

Nhớ một câu: Pearson hỏi \emph{có thẳng không}, Spearman hỏi \emph{có cùng chiều không}.

\subsection*{Hình dạng phân phối}

\textbf{Tính yếu vị:} một đỉnh / hai đỉnh / nhiều đỉnh / đều. Hai đỉnh thường là
dấu hiệu dữ liệu trộn hai nhóm --- nên tách ra thay vì báo cáo một trung bình vô
nghĩa nằm giữa.

\textbf{Độ lệch} nhìn theo \emph{đuôi dài}, không nhìn chỗ dồn cục. Lệch phải thì
$\bar x >$ trung vị; lệch trái thì $\bar x <$ trung vị. Nhớ bằng: \emph{trung bình
bị đuôi kéo}. Độ lệch quay lại ở kiểm định Jarque--Bera, mục~\ref{sec:kiemdinh}.

\textbf{Thống kê bền:} phân phối đối xứng thì báo cáo $\bar x$ và $s$; phân phối
lệch hoặc có ngoại lệ thì báo cáo \emph{trung vị và IQR}.

\textbf{Biến đổi log} xử lý dữ liệu lệch phải mạnh: kéo đuôi phải lại gần, biến
quan hệ nhân thành quan hệ cộng. Đây là lý do đồ án đo bằng
$Y = 100\log(P_1/P_0)$ chứ không lấy hiệu giá thô.

\subsection*{Bẫy}

\begin{itemize}
\item Mẫu số phương sai là $n-1$, không phải $n$. Chia $n$ (và dùng trung bình
      quần thể $\mu$) mới là $\sigma^2$.
\item $s^2$ mang đơn vị bình phương nên vô nghĩa khi diễn giải; $s$ mới cùng đơn
      vị với dữ liệu.
\item $r=0$ \emph{không} suy ra độc lập. $y=x^2$ với $x$ đối xứng quanh 0 cho
      $r=0$ nhưng $y$ hoàn toàn xác định bởi $x$.
\item Công thức rút gọn của Spearman \emph{chỉ đúng khi không trùng hạng}. Có
      trùng hạng thì phải tính Pearson trên hai cột hạng.
\item Phân biệt hai từ gần giống: \emph{đồng hạng} là cùng chiều (concordant),
      \emph{trùng hạng} là hai giá trị bằng nhau nên chia nhau một hạng (ties).
\end{itemize}

\subsection*{Ví dụ tính tay}

Dữ liệu $x = 1,2,3,4,5$ và $y = 2,4,5,4,9$. Ta có $\bar x = 3$, $\bar y = 4{,}8$,
$S_{xx}=10$, $S_{xy}=14$, $S_{yy}=26{,}8$.

\[
r = \frac{14}{\sqrt{10 \times 26{,}8}} = \frac{14}{16{,}37} = 0{,}855
\]

Hạng của $y$: hai giá trị $y=4$ chia nhau hạng 2 và 3 nên cùng nhận $2{,}5$, cho
hạng $(1;\,2{,}5;\,4;\,2{,}5;\,5)$. Khi đó $\sum d_i^2 = 3{,}5$ và

\[
\rho = 1 - \frac{6 \times 3{,}5}{5(25-1)} = 0{,}825
\]

Vì có trùng hạng nên đây là xấp xỉ. Giá trị chính xác (Pearson trên hai cột hạng)
là $8/\sqrt{95} = 0{,}8208$. Dấu hiệu nhận biết: tổng bình phương hai cột hạng
lệch nhau ($10$ so với $9{,}5$); không trùng hạng thì chúng luôn bằng nhau.

Kendall trên cùng dữ liệu: $C=8$, $D=1$, nên $\tau = 7/10 = 0{,}70$.

Ba hệ số cho ba con số khác nhau trên cùng một dữ liệu ($0{,}855$; $0{,}821$;
$0{,}70$) --- chúng đo ba thứ khác nhau nên không việc gì phải bằng nhau.

\subsection*{Dữ liệu định danh}

\textbf{Bảng ngẫu nhiên} tóm tắt hai biến định tính bằng tần suất. Tỷ lệ hàng và
tỷ lệ cột trả lời \emph{hai câu hỏi khác nhau}: hàng hỏi ``trong số người hút, bao
nhiêu phần trăm mắc bệnh'', cột hỏi ``trong số người mắc bệnh, bao nhiêu phần trăm
có hút''. Chọn sai chiều là diễn giải sai hoàn toàn --- cùng một bẫy đảo ngược xác
suất có điều kiện gặp lại ở p-value và ở độ nhạy/độ đặc hiệu.

Cột chồng chuẩn hoá là lựa chọn đúng khi hai nhóm chênh lệch cỡ; cột chồng thường
làm nhóm nhỏ trông như không có gì.

%=====================================================================
\section{Khoảng tin cậy (chương 5)}\label{sec:ktc}

\subsection*{Công thức}

Dạng chung: \textbf{ước lượng điểm} $\pm\ z^{*} \times SE$.

\begin{center}
\begin{tabular}{@{}ll@{}}
\toprule
\textbf{Tham số} & \textbf{Khoảng tin cậy} \\
\midrule
Tỉ lệ $p$ & $\hat p \pm z^{*}\sqrt{\hat p(1-\hat p)/n}$ \\[2pt]
Trung bình, biết $\sigma$ & $\bar x \pm z^{*}\,\sigma/\sqrt{n}$ \\[2pt]
Trung bình, không biết $\sigma$ \; (ngoài book) & $\bar x \pm t^{*}_{n-1}\,s/\sqrt{n}$ \\[2pt]
Hệ số góc & $b_1 \pm t^{*}_{n-2}\,SE_{b_1}$ \\
\bottomrule
\end{tabular}
\end{center}

\[
\hat{p} \sim \mathcal{N}\!\left(\mu = p,\ \ SE=\sqrt{\frac{p(1-p)}{n}}\right)
\qquad
n \ge \frac{(z^{*})^2\,p(1-p)}{ME^2}
\]

Chưa biết $p$ thì lấy $p=0{,}5$ (làm $p(1-p)$ lớn nhất, an toàn nhất) và luôn làm
tròn lên.

\textbf{Điều kiện dùng xấp xỉ chuẩn:} độc lập, và ít nhất 10 thành công \emph{và}
10 thất bại kỳ vọng ($np \ge 10$, $n(1-p) \ge 10$).

\subsection*{Bẫy}

\begin{itemize}
\item Dựng khoảng tin cậy thì $SE$ dùng $\hat p$; kiểm định thì $SE$ dùng $p_0$.
      Đây là khác biệt giữa mục 5.2 và 5.3 của book.
\item \emph{Sai số mẫu} giảm khi tăng $n$; \emph{độ chệch} thì không. Mẫu bị chệch
      mà tăng $n$ lên 100.000 chỉ cho một khoảng rất hẹp bao quanh con số sai.
\item $ME \propto 1/\sqrt{n}$: muốn giảm sai số một nửa phải tăng mẫu \emph{gấp bốn}.
\item Đề cho $s$ thì dùng $t$ với $df=n-1$; cho $\sigma$ thì dùng $z$.
\end{itemize}

\subsection*{Đọc kết quả cho đúng}

Với khoảng tin cậy 95\% bằng $(0{,}64;\ 0{,}70)$:

\begin{itemize}
\item \textbf{Sai:} ``95\% người dùng nằm trong khoảng 64--70\%'' --- khoảng tin
      cậy nói về \emph{tham số}, không nói về cá thể.
\item \textbf{Sai:} ``xác suất $p$ nằm trong khoảng là 95\%'' --- $p$ là hằng số,
      cái ngẫu nhiên là \emph{khoảng}.
\item \textbf{Đúng:} ``nếu lặp lại việc lấy mẫu và dựng khoảng, khoảng 95\% số
      khoảng sẽ chứa $p$''. Viết ngắn được chấp nhận: ``ta tin cậy 95\% rằng $p$
      nằm trong $(0{,}64;\ 0{,}70)$''.
\end{itemize}

\subsection*{Ví dụ}

\textbf{(a) Tỉ lệ.} $n=850$, $\hat p = 0{,}67$, độ tin cậy 95\%.
$SE = \sqrt{0{,}67 \times 0{,}33/850} = 0{,}0161$, biên sai số
$1{,}96 \times 0{,}0161 = 0{,}032$, khoảng $(0{,}64;\ 0{,}70)$.

\textbf{(b) Trung bình, chưa biết $\sigma$.} $n=25$, $\bar x = 52$, $s=8$.
$df=24$, $t^{*}=2{,}064$, $SE = 8/5 = 1{,}6$, biên sai số $3{,}30$, khoảng
$(48{,}70;\ 55{,}30)$. Lỡ dùng $z^{*}=1{,}96$ sẽ ra $(48{,}86;\ 55{,}14)$ ---
hẹp hơn, tức tự tin quá mức, đúng cái lỗi mà phân phối $t$ sinh ra để sửa.

\textbf{(c) Cỡ mẫu.} Muốn $ME \le 3\%$ ở 95\%, chưa biết $p$:
$n \ge 1{,}96^2 \times 0{,}25/0{,}03^2 = 1067{,}1$, lấy $n=1068$. Đây là lý do các
cuộc thăm dò dư luận toàn thấy $n \approx 1000$, sai số $\pm 3\%$.

\subsection*{Cầu nối sang kiểm định}

$p_0$ \emph{nằm ngoài} khoảng tin cậy $(1-\alpha)$ khi và chỉ khi bác bỏ
$H_0: p=p_0$ ở mức $\alpha$ (hai phía). Hai thứ là một phép tính nhìn từ hai phía.

%=====================================================================
\section{Kiểm định giả thuyết (chương 6)}\label{sec:kiemdinh}

Chương nặng nhất. Mục 6.2 của book là một bảng khoảng tám loại kiểm định; đề gần
như chắc chắn rút từ đó. Thuộc bảng chọn công thức trước, hiểu sau.

\subsection*{Bốn bước}

Phát biểu $H_0$ và $H_a$ $\to$ chọn thống kê kiểm định và phân phối của nó
\emph{khi $H_0$ đúng} $\to$ xác định miền bác bỏ theo $\alpha$ $\to$ bác bỏ nếu
giá trị tính được rơi vào miền đó. Lối hiện đại tương đương: tính thống kê, tính
p-value, bác bỏ nếu \textbf{p-value $< \alpha$}.

\subsection*{Đặt giả thuyết cho đúng chiều}

\begin{center}
\begin{tabular}{@{}lll@{}}
\toprule
\textbf{Đề nói} & \textbf{$H_a$} & \textbf{Miền bác bỏ} \\
\midrule
có khác không, có thay đổi không & $\mu \ne \mu_0$ & hai đuôi, $|z| > z_{\alpha/2}$ \\
có lớn hơn không, có tăng không & $\mu > \mu_0$ & đuôi phải, $z > z_\alpha$ \\
có nhỏ hơn không, có giảm không & $\mu < \mu_0$ & đuôi trái, $z < -z_\alpha$ \\
\bottomrule
\end{tabular}
\end{center}

$H_0$ luôn mang dấu bằng; điều muốn chứng minh nằm ở $H_a$. Chọn chiều
\emph{trước khi nhìn dữ liệu}. Với $\alpha = 0{,}05$: một phía dùng $1{,}645$, hai
phía dùng $1{,}96$.

\subsection*{Hai loại sai lầm}

\begin{center}
\begin{tabular}{@{}lll@{}}
\toprule
 & \textbf{Không bác bỏ $H_0$} & \textbf{Bác bỏ $H_0$} \\
\midrule
$H_0$ đúng & đúng & \textbf{sai lầm loại I} ($\alpha$), dương tính giả \\
$H_0$ sai & \textbf{sai lầm loại II} ($\beta$), âm tính giả & đúng \\
\bottomrule
\end{tabular}
\end{center}

Lực kiểm định $= 1-\beta$. Giữ nguyên $n$ mà giảm $\alpha$ thì $\beta$ tăng; muốn
giảm cả hai chỉ còn cách tăng $n$.

\subsection*{Bảng chọn công thức}

\begin{center}
\begin{tabular}{@{}lll@{}}
\toprule
\textbf{Bài toán} & \textbf{Thống kê} & \textbf{$df$} \\
\midrule
1 trung bình, biết $\sigma$ &
  $z = \dfrac{\bar x - \mu_0}{\sigma/\sqrt n}$ & --- \\[10pt]
1 trung bình, không biết $\sigma$ &
  $t = \dfrac{\bar x - \mu_0}{s/\sqrt n}$ & $n-1$ \\[10pt]
1 phương sai &
  $q = \dfrac{(n-1)s^2}{\sigma_0^2}$ & $n-1$ \\[10pt]
1 tỉ lệ &
  $z = \dfrac{\hat p - p_0}{\sqrt{p_0(1-p_0)/n}}$ & --- \\[10pt]
2 trung bình, phương sai bằng nhau &
  $t = \dfrac{\bar x_1 - \bar x_2}{\sqrt{s_p^2(1/n_1+1/n_2)}}$ & $n_1+n_2-2$ \\[10pt]
2 trung bình, $n_1,n_2 > 30$ &
  $z = \dfrac{\bar x_1 - \bar x_2}{\sqrt{s_1^2/n_1+s_2^2/n_2}}$ & --- \\[10pt]
2 phương sai &
  $f = \dfrac{s_1^2}{s_2^2}$ & $n_1-1,\ n_2-1$ \\[10pt]
2 tỉ lệ &
  $z = \dfrac{\hat p_1-\hat p_2}{\sqrt{\bar p(1-\bar p)(1/n_1+1/n_2)}}$ & --- \\[10pt]
Độc lập của hai biến định tính &
  $q = \sum\dfrac{(O-E)^2}{E}$ & $(h-1)(k-1)$ \\[10pt]
Tính chuẩn (Jarque--Bera) &
  $q = n\!\left(\dfrac{\text{skew}^2}{6}+\dfrac{\text{kurt}^2}{24}\right)$ & 2 \\[10pt]
Phù hợp phân phối (Kolmogorov--Smirnov) &
  $d = \max|F_o - F_m|$ & --- \\
\bottomrule
\end{tabular}
\end{center}

\[
s_p^2 = \frac{(n_1-1)s_1^2+(n_2-1)s_2^2}{n_1+n_2-2}
\qquad
\bar p = \frac{n_1\hat p_1+n_2\hat p_2}{n_1+n_2}
\qquad
E_{ij} = \frac{R_i \times C_j}{n}
\]

với $R_i$ là tổng hàng $i$ và $C_j$ là tổng cột $j$.

Phân phối $\chi^2$ không đối xứng nên hai đuôi phải tra hai giá trị khác nhau,
không lấy trị tuyệt đối được. Kiểm định hai phương sai thường chạy \emph{trước}
kiểm định hai trung bình, để quyết định có được gộp phương sai không.

\subsection*{Bài giải mẫu}

\textbf{Bài 1 --- một trung bình.} Giám đốc nói lương trung bình là 380; mẫu
$n=36$ cho $\bar x = 350$, $s=40$; $\alpha = 0{,}05$.
Có $s$ nên dùng $t$ với $df=35$. $SE = 40/6 = 6{,}667$ và
$t = (350-380)/6{,}667 = -4{,}50$. Miền bác bỏ $|t| > t_{0{,}025;\,35} = 2{,}030$.
Vì $4{,}50 > 2{,}030$ nên \textbf{bác bỏ} $H_0$ (p-value $\approx 0{,}00007$).

\textbf{Bài 2 --- hai trung bình, mẫu lớn.} Đường huyết trước:
$n_1=50$, $\bar x = 60$, $s_1=7$. Sau: $n_2=40$, $\bar y = 52$, $s_2=9{,}2$. Hỏi
có giảm không, tức một phía.
$SE = \sqrt{49/50 + 84{,}64/40} = \sqrt{3{,}096} = 1{,}760$ và
$z = 8/1{,}760 = 4{,}55 > 1{,}645$ nên \textbf{bác bỏ}. Chú ý: hai nhóm cỡ khác
nhau (50 và 40) nên đây \emph{không} phải dữ liệu ghép cặp dù đề nói ``trước và sau''.

\textbf{Bài 3 --- một tỉ lệ.} Phế phẩm trước đây 5\%; kiểm 800 sản phẩm thấy 24
phế phẩm; $\alpha=0{,}01$; đuôi trái.
$\hat p = 0{,}03$, $SE = \sqrt{0{,}05 \times 0{,}95/800} = 0{,}007706$,
$z = -0{,}02/0{,}007706 = -2{,}60 < -2{,}326$ nên \textbf{bác bỏ}
(p-value $=0{,}0047$). Nếu đề để $\alpha = 0{,}001$ thì ngưỡng là $-3{,}09$ và
\emph{không} bác bỏ được --- cùng dữ liệu, kết luận đổi theo $\alpha$.

\textbf{Bài 4 --- hai tỉ lệ.} Dùng DDT: 90 người, 10 nhiễm. Không dùng: 100 người,
26 nhiễm.
$\hat p_1 = 0{,}1111$, $\hat p_2 = 0{,}26$, $\bar p = 36/190 = 0{,}1895$,
$SE = \sqrt{0{,}1895 \times 0{,}8105 \times (1/90+1/100)} = 0{,}05694$,
$z = -0{,}14889/0{,}05694 = -2{,}615$. Vì $2{,}615 > 1{,}96$ nên \textbf{bác bỏ}.
Tử số dùng $\hat p_1, \hat p_2$ riêng, mẫu số dùng $\bar p$ gộp --- vì mẫu số
được tính \emph{dưới giả định $H_0$ đúng}, mà $H_0$ nói hai tỉ lệ bằng nhau.

\bay Giữ tử số ở 5 chữ số là có chủ ý. Lấy tử số làm tròn thành $-0{,}1489$ rồi
chia sẽ ra $2{,}61503$, tức $2{,}62$; tính đủ chữ số mới ra $2{,}61487$, tức
$2{,}615$. Cùng cái bẫy làm tròn ở mục hồi quy, chỉ khác là ở đây nó đủ sức đổi
chữ số cuối của đáp án.

\textbf{Bài 5 --- chi bình phương độc lập.} Bảng quan sát: hút thuốc có 30 mắc
bệnh và 70 không (tổng 100); không hút có 20 mắc và 130 không (tổng 150). Tổng
cột là 50 và 200, tổng cộng $n=250$.
Kỳ vọng lần lượt là $20$; $80$; $30$; $120$ (đều $\ge 5$ nên dùng được).
\[
q = \frac{10^2}{20}+\frac{10^2}{80}+\frac{10^2}{30}+\frac{10^2}{120}
  = 5 + 1{,}25 + 3{,}333 + 0{,}833 = 10{,}42
\]
$df = 1$, ngưỡng $\chi^2_{0{,}05;1} = 3{,}841$, nên \textbf{bác bỏ}: hai biến
\emph{không độc lập}. Viết kết luận là ``không độc lập'', \emph{đừng} viết ``hút
thuốc gây bệnh phổi'' --- chi bình phương chỉ phát hiện liên hệ.

\subsection*{Bẫy}

\begin{itemize}
\item p-value là $P(\mathcal{D} \mid H_0)$, \emph{không} phải xác suất $H_0$ đúng.
      Đảo ngược hai vế của xác suất có điều kiện là lỗi bị trừ điểm nặng nhất.
\item Quên chia đôi $\alpha$ ở kiểm định hai phía.
\item Dùng $\hat p$ ở mẫu số khi kiểm định một tỉ lệ; phải là $p_0$.
\item Quên kiểm điều kiện ($np \ge 10$, mọi $E_{ij} \ge 5$) --- nhiều đề cho điểm
      riêng cho bước này.
\item Viết ``chấp nhận $H_0$''. Phải viết \emph{chưa đủ bằng chứng để bác bỏ $H_0$}.
\item Nhầm ý nghĩa thống kê với ý nghĩa thực tiễn: $n$ đủ lớn thì chênh lệch bé
      xíu cũng có $p<0{,}05$.
\item Đa kiểm định: chạy 20 kiểm định độc lập ở $\alpha = 0{,}05$ với mọi $H_0$
      đều đúng thì kỳ vọng vẫn có $20 \times 0{,}05 = 1$ cái ``có ý nghĩa''.
\end{itemize}

%=====================================================================
\section{Mô hình đồ thị (chương 7)}\label{sec:dothi}

Chương rẻ điểm nhất --- không phải tính toán, chỉ cần đọc đúng đồ thị.

\subsection*{Mạng Bayes}

\[
p(x_1,\dots,x_n) = \prod_{i=1}^{n} p\bigl(x_i \mid pa(x_i)\bigr)
\]

Đọc ngược một phân rã ra đồ thị: vế điều kiện chính là danh sách nút cha.
\emph{Giả định cạnh nhân quả:} mọi nút cha là nguyên nhân trực tiếp của tất cả nút
con của nó. Đây là cái biến một mạng Bayes thuần xác suất thành đồ thị nhân quả,
và nó đến từ \emph{kiến thức lĩnh vực} chứ không từ dữ liệu.

\subsection*{Ba cấu trúc ba nút --- trái tim của chương}

\begin{center}
\begin{tabular}{@{}llll@{}}
\toprule
\textbf{Cấu trúc} & \textbf{Vai trò của $Z$} & \textbf{Chưa điều kiện $Z$} & \textbf{Điều kiện $Z$} \\
\midrule
$X \to Z \to Y$ (chuỗi) & trung gian & mở & \textbf{chặn} \\
$X \leftarrow Z \to Y$ (rẽ nhánh) & gây nhiễu & mở & \textbf{chặn} \\
$X \to Z \leftarrow Y$ (hội tụ) & hội tụ & \textbf{chặn} & \textbf{mở} \\
\bottomrule
\end{tabular}
\end{center}

Collider đi \emph{ngược} lại với hai cái kia --- đó là toàn bộ độ khó của chương.

\textbf{Quy tắc chặn đường.} Một đường bị chặn nếu chứa một biến \emph{không hội
tụ được} điều kiện hoá, hoặc một biến \emph{hội tụ không được} điều kiện hoá (và
không hậu duệ nào của nó được điều kiện hoá). Hai biến độc lập có điều kiện khi
\emph{mọi} đường giữa chúng đều bị chặn --- chỉ cần một đường mở là còn liên hệ.

\subsection*{Sáu vai trò của biến}

Với $T$ là tác nhân và $Y$ là kết quả:

\begin{center}
\begin{tabular}{@{}llp{7.2cm}@{}}
\toprule
\textbf{Vai trò} & \textbf{Kiểm soát?} & \textbf{Vì sao} \\
\midrule
Gây nhiễu $T \leftarrow C \to Y$ & \textbf{phải} & không kiểm soát thì liên hệ bị pha tạp \\
Trung gian $T \to M \to Y$ & không & kiểm soát là chặn mất một phần tác động cần đo \\
Hội tụ $T \to K \leftarrow Y$ & tuyệt đối không & kiểm soát nó \emph{tạo ra} liên hệ giả \\
Hậu duệ của hội tụ & không & mở đường hội tụ một phần \\
Tiền định $P \to Y$ & nên & không bắt buộc, nhưng giảm phương sai \\
Công cụ $I \to T$ & không & làm \emph{tăng} phương sai \\
\bottomrule
\end{tabular}
\end{center}

Một câu nhớ cả bảng: \textbf{chặn cửa sau, đừng chạm cửa trước, đừng bao giờ mở
cửa hội tụ.}

\subsection*{Ví dụ 1 --- đọc phân rã ra đồ thị}

Cho $p(l,g,i,d,s) = p(l\mid g)\,p(g\mid i,d)\,p(i)\,p(d)\,p(s\mid i)$, tức các
cạnh $i \to g$, $d \to g$, $i \to s$, $g \to l$.

\begin{itemize}
\item $s$ và $l$ \emph{không} độc lập: đường $s \leftarrow i \to g \to l$ có $i$ là
      nút rẽ nhánh và $g$ là nút chuỗi, cả hai chưa điều kiện nên đường mở.
\item $s \perp l \mid g$: điều kiện hoá $g$ (nút chuỗi) làm chặn đường duy nhất đó.
\item $s$ và $d$ \emph{độc lập}: đường $s \leftarrow i \to g \leftarrow d$ có $g$
      là collider chưa điều kiện nên chặn sẵn.
\item $s$ và $d$ \emph{không còn} độc lập khi biết $g$: điều kiện hoá đúng collider
      thì mở đường. Có lý giải rõ --- biết một sinh viên điểm thấp, lại biết môn
      đó dễ, ta suy ra bạn ấy có lẽ không thông minh, mà thế thì điểm SAT cũng
      thấp. Điều kiện hoá $g$ đã bơm thông tin chảy giữa hai thứ vốn không liên quan.
\end{itemize}

\subsection*{Ví dụ 2 --- chọn tập kiểm soát}

Cạnh: $C \to T$, $C \to Y$, $T \to M$, $M \to Y$, $T \to K$, $Y \to K$.
Liệt kê mọi đường từ $T$ tới $Y$:

\begin{center}
\begin{tabular}{@{}lll@{}}
\toprule
\textbf{Đường} & \textbf{Nút giữa} & \textbf{Trạng thái} \\
\midrule
$T \to M \to Y$ & $M$ chuỗi & mở --- \emph{đây là tác động thật} \\
$T \leftarrow C \to Y$ & $C$ rẽ nhánh & mở --- cửa sau, \emph{phải chặn} \\
$T \to K \leftarrow Y$ & $K$ hội tụ & chặn sẵn --- \emph{đừng đụng vào} \\
\bottomrule
\end{tabular}
\end{center}

Tập kiểm soát đúng là $\{C\}$. Ba lựa chọn sai: để trống thì còn nhiễu bởi $C$;
lấy $\{C,M\}$ thì chặn luôn tác động cần đo, chỉ còn tác động trực tiếp; lấy
$\{C,K\}$ thì tự tạo ra liên hệ giả.

\textbf{Bài học:} thêm biến kiểm soát không phải lúc nào cũng tốt hơn. ``Cho hết
vào mô hình cho chắc'' là sai lầm nghiêm trọng, và đây chính là điều chương 7 dạy.

\subsection*{Liên hệ đồ án}

Trong đồ án thuế GTGT, $D$ (thuế suất cửa hàng thực áp) nằm \emph{giữa} $Z$ (đủ
điều kiện giảm thuế) và $Y$, tức là \emph{biến trung gian}. Đó là lý do đồ án ước
lượng theo $Z$ chứ không hồi quy theo $D$ --- điều kiện hoá theo $D$ sẽ chặn mất
chính con đường cần đo. Câu này rất dễ bị hỏi khi vấn đáp.

%=====================================================================
\section{Phân tích nhân quả (chương 8)}

\subsection*{Ba cấp độ}

Liên hệ (thấy gì, học được từ dữ liệu quan sát thụ động) $\to$ can thiệp (nếu ta
làm thì sao, cần mô hình nhân quả hoặc thí nghiệm) $\to$ phản thực (giá như đã
không làm thì sao, cần thêm giả định về từng cá thể). Học máy thông thường nằm ở
\emph{cấp một} --- dự đoán giỏi không đồng nghĩa với hiểu nhân quả.

\subsection*{Kết quả tiềm năng}

Với cá thể $i$: $Y_i(1)$ là kết quả giá như được can thiệp, $Y_i(0)$ là giá như
không. Tác động nhân quả cá thể là $Y_i(1)-Y_i(0)$.

\textbf{Bài toán nền tảng:} ta chỉ quan sát được \emph{một} trong hai. Tác động
nhân quả cá thể là không quan sát được --- không phải vì thiếu dữ liệu mà vì
logic. Nên chuyển sang đại lượng trung bình:

\[
ATE = \mathbb{E}\bigl[Y_i(1)-Y_i(0)\bigr] = \mathbb{E}[Y(1)] - \mathbb{E}[Y(0)]
\]

\subsection*{Bất đẳng thức trung tâm}

\[
\underbrace{\mathbb{E}[Y(1)] - \mathbb{E}[Y(0)]}_{\text{nhân quả}}
\ \neq\
\underbrace{\mathbb{E}[Y\mid T{=}1] - \mathbb{E}[Y\mid T{=}0]}_{\text{liên hệ}}
\]

Vế trái so cùng một nhóm người ở hai kịch bản, không quan sát trực tiếp được. Vế
phải so hai nhóm người khác nhau, tính thẳng từ dữ liệu. Chênh lệch giữa chúng:

\[
\underbrace{\mathbb{E}[Y\mid T{=}1] - \mathbb{E}[Y\mid T{=}0]}_{\text{(1) đo được}}
= \underbrace{ATT}_{\text{(2) cái cần}}
+ \underbrace{\bigl(\mathbb{E}[Y(0)\mid T{=}1] - \mathbb{E}[Y(0)\mid T{=}0]\bigr)}_{\text{(3) chệch chọn lọc}}
\]

Số hạng (3) hỏi: \emph{hai nhóm có khác nhau ngay cả khi không ai được can thiệp
không?} Ví dụ người đi gym khoẻ hơn người không đi, nhưng người vốn khoẻ mới hay
đi gym, nên hiệu quả gym bị thổi phồng.

\textbf{RCT giải quyết bằng cách nào.} Phân bổ ngẫu nhiên làm $T$ độc lập với mọi
kết quả tiềm năng, nên (3) bằng 0 và liên hệ \emph{trở thành} nhân quả. Điểm mấu
chốt: ngẫu nhiên hoá cân bằng cả những biến gây nhiễu \emph{ta không biết và không
đo được}. Không phương pháp thống kê nào trên dữ liệu quan sát làm được điều đó.

\subsection*{Hai nghịch lý}

\textbf{Simpson --- bẫy gây nhiễu.} Xu hướng trong từng nhóm con đảo ngược khi gộp.
Ví dụ sỏi thận: cách A khỏi $81/87 = 93\%$ với sỏi nhỏ và $192/263 = 73\%$ với sỏi
lớn; cách B khỏi $234/270 = 87\%$ và $55/80 = 69\%$. A thắng trong \emph{cả hai}
nhóm, nhưng gộp lại thì A được $273/350 = 78\%$ còn B được $289/350 = 83\%$ --- B
thắng. Nguyên nhân: kích thước sỏi vừa ảnh hưởng việc chọn cách chữa vừa ảnh hưởng
khả năng khỏi, tức là \emph{biến gây nhiễu}, nên phải đọc con số \emph{theo nhóm con}.

\textbf{Berkson --- bẫy hội tụ.} Một nhóm con thể hiện liên hệ mà quần thể rộng
không có. Trong giới nổi tiếng, tài năng và ngoại hình tương quan âm giả tạo, vì
muốn nổi tiếng thì cần \emph{hoặc} tài \emph{hoặc} đẹp. Chọn mẫu chỉ gồm người nổi
tiếng chính là điều kiện hoá đúng nút hội tụ.

\subsection*{Tương quan không phải nhân quả}

Thấy $X$ và $Y$ tương quan thì có bốn khả năng: $X \to Y$; $Y \to X$; có biến gây
nhiễu chung $X \leftarrow C \to Y$; hoặc ngẫu nhiên / chệch chọn mẫu. Dữ liệu quan
sát không tự phân biệt được --- cần thêm chiều thời gian, kiến thức lĩnh vực, hoặc
can thiệp.

\subsection*{Ngoài book --- cho vấn đáp đồ án}

\textbf{Tiêu chuẩn cửa sau:} tập $C$ chặn mọi đường cửa sau từ $T$ tới $Y$ và không
chứa hậu duệ của $T$; khi đó
$P\bigl(Y \mid do(T{=}t)\bigr) = \sum_c P(Y \mid T{=}t, C{=}c)\,P(C{=}c)$.

\textbf{g-computation:} khớp $\hat m_0(X)$ \emph{chỉ trên nhóm chứng}, rồi dùng
chính mô hình đó dự đoán kết quả phản thực cho nhóm can thiệp:
$\widehat{ATT} = \frac{1}{n_1}\sum_{i:T_i=1}\bigl(Y_i - \hat m_0(X_i)\bigr)$.

%=====================================================================
\section{Hồi quy tuyến tính (chương 9)}

Chương được điểm nhanh nhất --- bài tập gần như chỉ có một dạng.

\subsection*{Công thức}

\[
\hat y = \beta_0 + \beta_1 x
\qquad
e_i = y_i - \hat y_i
\qquad
b_1 = \frac{s_y}{s_x}R = \frac{S_{xy}}{S_{xx}}
\qquad
b_0 = \bar y - b_1 \bar x
\]

\[
R^2 = \frac{SS_{\text{Model}}}{SS_{\text{Total}}} = 1 - \frac{SS_{\text{Error}}}{SS_{\text{Total}}}
\qquad
s_e = \sqrt{\frac{SS_{\text{Error}}}{n-2}}
\qquad
SE_{b_1} = \frac{s_e}{\sqrt{S_{xx}}}
\]

\[
T = \frac{b_1 - \beta_1^{0}}{SE_{b_1}}
\qquad
df = n-2
\qquad
CI:\ b_1 \pm t^{*}_{n-2}\,SE_{b_1}
\]

Bình phương tối thiểu chọn hệ số để $\sum e_i^2$ nhỏ nhất --- không phải
$\sum|e_i|$, cũng không phải $\sum e_i$ (tổng này luôn bằng 0). Trong hồi quy đơn,
$R^2$ đúng bằng bình phương hệ số tương quan.

\textbf{Ba tính chất dùng để tự kiểm tra:} $b_1$ luôn cùng dấu với $R$; đường hồi
quy luôn đi qua $(\bar x, \bar y)$; tổng phần dư luôn bằng 0.

\subsection*{Điều kiện và ngoại lệ}

Ba điều kiện: quan hệ tuyến tính, phần dư gần chuẩn, phương sai không đổi (cộng
thêm điều kiện ngầm là quan sát độc lập). Ba loại ngoại lệ: xa theo chiều dọc
(ngoại lệ phần dư), xa theo chiều ngang (đòn bẩy cao), và \emph{có ảnh hưởng}
(bỏ ra thì hệ số góc đổi đáng kể) --- loại thứ ba nguy hiểm nhất.

\subsection*{Ví dụ 1 --- từ tóm tắt thống kê}

$\bar x = 86{,}01$ (\% tốt nghiệp), $\bar y = 11{,}35$ (\% nghèo), $s_x = 3{,}73$,
$s_y = 3{,}1$, $R = -0{,}75$.

$b_1 = (3{,}1/3{,}73)(-0{,}75) = -0{,}62$ và
$b_0 = 11{,}35 + 0{,}62 \times 86{,}01 = 64{,}68$, cho $\hat y = 64{,}68 - 0{,}62x$.
Dự đoán tại $x = 85{,}1$ được $11{,}92\%$. Hệ số xác định $R^2 = 0{,}5625$.

Diễn giải: mỗi điểm phần trăm tốt nghiệp tăng thêm \emph{đi kèm} mức nghèo thấp
hơn $0{,}62$ điểm phần trăm. Viết ``đi kèm'', đừng viết ``làm giảm'' --- đây là dữ
liệu quan sát. Hệ số chặn $64{,}68$ ứng với $x=0$, nằm ngoài xa dải dữ liệu nên
\emph{vô nghĩa}; luôn nhắc một câu về ngoại suy khi làm bài.

\bay Giá trị chưa làm tròn là $b_1 = -0{,}62332$; giữ nguyên nó thì $b_0 = 64{,}96$
chứ không phải $64{,}68$, lệch $0{,}28$. Vì $b_1$ bị nhân với $\bar x = 86$, sai số
làm tròn ở chữ số thứ ba bị khuếch đại gần 90 lần. Giữ đủ chữ số cho tới bước cuối.

\subsection*{Ví dụ 2 --- tính tay đầy đủ}

Dùng lại $x = 1,\dots,5$ và $y = 2,4,5,4,9$ ở mục~1: $S_{xy}=14$, $S_{xx}=10$,
$S_{yy}=26{,}8$.

$b_1 = 14/10 = 1{,}4$ và $b_0 = 4{,}8 - 4{,}2 = 0{,}6$, cho $\hat y = 0{,}6+1{,}4x$.
Giá trị khớp là $2{,}0$; $3{,}4$; $4{,}8$; $6{,}2$; $7{,}6$, nên phần dư là
$0$; $0{,}6$; $0{,}2$; $-2{,}2$; $1{,}4$ --- cộng lại đúng bằng 0.

$SS_{\text{Error}} = 7{,}20$ nên $R^2 = 1 - 7{,}20/26{,}80 = 0{,}731$
(kiểm chéo: $0{,}855^2 = 0{,}731$).
$s_e = \sqrt{7{,}20/3} = 1{,}549$ và $SE_{b_1} = 1{,}549/\sqrt{10} = 0{,}490$, cho
$T = 1{,}4/0{,}490 = 2{,}86$ với $df=3$.

Ngưỡng $t^{*}_{0{,}025;3} = 3{,}182 > 2{,}86$ nên \emph{không} bác bỏ
$H_0: \beta_1 = 0$. Khoảng tin cậy
$1{,}4 \pm 3{,}182 \times 0{,}490 = (-0{,}16;\ 2{,}96)$ chứa 0, nhất quán.

\textbf{Bài học quan trọng nhất.} $R^2 = 73\%$ trông rất đẹp nhưng hệ số góc
\emph{không có ý nghĩa thống kê}, vì $n=5$ quá nhỏ. $R^2$ đo \emph{khớp bao
nhiêu}, còn $p$ đo \emph{có chắc không phải may rủi không} --- hai chuyện khác nhau.

\subsection*{Ví dụ 3 --- đọc bảng output}

Nghiên cứu IQ song sinh, $n=27$: hệ số góc $0{,}9014$ với sai số chuẩn $0{,}0963$.

Kiểm chứng cột $t$: $0{,}9014/0{,}0963 = 9{,}36$ --- cột $t$ luôn là ước lượng chia
sai số chuẩn, tự tính lại để chắc mình đọc đúng cột. Với $df = 25$ và
$t^{*} = 2{,}06$, khoảng tin cậy là $(0{,}70;\ 1{,}10)$. Khoảng này \emph{chứa 1},
nên không bác bỏ được giả thuyết $\beta_1 = 1$, tức quan hệ có thể là một đổi một.

\bay Đề hỏi kiểm định $\beta_1 = 1$ thì tử số là $b_1 - 1$. Cột $t$ trong bảng
output \emph{luôn} ứng với null bằng 0, phải tự tính lại.

%=====================================================================
\section{Hồi quy đa biến và logistic (chương 10)}

\subsection*{Hồi quy đa biến}

\[
\hat y = \beta_0 + \beta_1 x_1 + \dots + \beta_p x_p
\qquad
R^2_{adj} = 1 - \left(\frac{SS_{\text{Error}}}{SS_{\text{Total}}} \times \frac{n-1}{n-p-1}\right)
\]

Diễn giải $\beta_j$: khi $x_j$ tăng 1 đơn vị \emph{và giữ nguyên mọi biến còn lại}.
Bỏ cụm đó là mất điểm diễn giải. Cùng một biến, hệ số trong mô hình đơn và mô hình
đa biến có thể khác nhau, thậm chí đổi dấu.

$R^2$ \emph{không bao giờ giảm} khi thêm biến, kể cả biến rác, nên vô dụng để so
mô hình khác số biến. $R^2_{adj}$ phạt theo $p$ nên \emph{có thể} giảm --- đó là
tính năng. So mô hình khác số biến thì dùng $R^2_{adj}$; báo cáo mức khớp của một
mô hình thì dùng $R^2$.

\vd $SS_{\text{Total}} = 480{,}25$, $SS_{\text{Error}} = 347{,}68$, $n=51$, $p=1$.
Khi đó $R^2 = 132{,}57/480{,}25 = 0{,}276$ và
$R^2_{adj} = 1 - (347{,}68/480{,}25)(50/49) = 0{,}261$.

\subsection*{Biến định tính}

Biến $k$ mức cần $k-1$ dummy; mức bị bỏ ra là \emph{mức tham chiếu} và hệ số chặn
đại diện cho nó. Cho cả $k$ dummy thì tổng của chúng luôn bằng 1, trùng với hệ số
chặn, ma trận suy biến --- đây là \emph{bẫy biến giả}.

\vd Với bốn vùng và northeast làm tham chiếu,
$\hat y = 9{,}50 + 0{,}03\,\text{midwest} + 1{,}79\,\text{west} + 4{,}16\,\text{south}$.
Miền Nam nghèo hơn Đông Bắc $4{,}16$ điểm phần trăm. \emph{Mọi hệ số dummy đều so
với mức tham chiếu}, không so với trung bình chung; muốn so midwest với south thì
lấy hiệu $4{,}16 - 0{,}03 = 4{,}13$.

\vd Dummy cộng biến liên tục: với $V$ là thể tích và $B=1$ nếu bìa mềm,
$\hat y = 197{,}96 + 0{,}72V - 184{,}05B$ tách thành hai đường \emph{song song} ---
bìa cứng $\hat y = 197{,}96 + 0{,}72V$ và bìa mềm $\hat y = 13{,}91 + 0{,}72V$.
Muốn hai đường khác hệ số góc thì phải thêm số hạng tương tác $V \times B$.

\subsection*{Lựa chọn mô hình và đa cộng tuyến}

\emph{Loại lùi:} bắt đầu từ mô hình đầy đủ, bỏ biến có p-value cao nhất, khớp lại,
lặp. \emph{Chọn tiến:} khớp hồi quy đơn cho từng biến, thêm biến có p-value nhỏ
nhất, lặp. Hai cách không nhất thiết cho cùng kết quả.

Đa cộng tuyến: thêm biến làm $R^2$ nhích lên nhưng $R^2_{adj}$ đứng yên nghĩa là
biến mới gần như không mang thông tin mới. Hậu quả là sai số chuẩn phình to và hệ
số bất ổn định. Ưu tiên mô hình gọn.

\subsection*{Hồi quy logistic}

\[
\text{odds} = \frac{p}{1-p}
\qquad
\text{logit}(p) = \log\frac{p}{1-p}
\qquad
p = \frac{e^{\eta}}{1+e^{\eta}},\ \ \eta = \beta_0 + \beta_1 x_1 + \dots
\]

Logit kéo $p \in [0,1]$ ra thành $(-\infty,+\infty)$ --- nhờ vậy mới đặt được mô
hình tuyến tính lên. Hệ số $\beta_j$ là thay đổi \emph{log-odds}; $e^{\beta_j}$ là
\emph{tỉ số odds}. Mốc ``không tác dụng'' là $e^{\beta}=1$, không phải 0.

\textbf{Ba bước bắt buộc:} logit $\to$ odds ($e^{\text{logit}}$) $\to$ xác suất
($\text{odds}/(1+\text{odds})$). Nhảy tắt là sai.

\vd Đoàn Donner, $\text{logit} = 1{,}8185 - 0{,}0665 \times \text{tuổi}$.
Người 25 tuổi: logit $= 0{,}1560$, odds $= e^{0{,}156} = 1{,}169$, xác suất
$= 1{,}169/2{,}169 = 0{,}539$. Người 50 tuổi: logit $= -1{,}5065$, odds $=0{,}2217$,
xác suất $= 0{,}181$. Tỉ số odds theo tuổi $e^{-0{,}0665} = 0{,}936$, tức mỗi năm
tuổi odds sống sót giảm khoảng $6{,}4\%$.

Với hệ số giới tính $1{,}5973$ thì $e^{1{,}5973} = 4{,}94$: giữ nguyên tuổi, odds
sống sót của nữ cao gấp khoảng 4,9 lần nam. Đừng viết ``xác suất cao gấp 5 lần'' ---
gấp 5 lần là \emph{odds}. Tỉ số odds cũng không phải tỉ số nguy cơ; chúng chỉ gần
nhau khi biến cố hiếm.

\subsection*{Độ nhạy, độ đặc hiệu, ngưỡng}

\[
\text{Sensitivity} = \frac{TP}{TP+FN}
\qquad
\text{Specificity} = \frac{TN}{FP+TN}
\]

Đổi ngưỡng thì hai chỉ số đi \emph{ngược chiều} nhau. Không có ngưỡng đúng về mặt
toán học --- chọn ngưỡng là quyết định kinh doanh, tuỳ giá của mỗi loại sai.

\vd Lọc thư rác ở ngưỡng $0{,}75$ với $TP=27$, $TN=3545$, $FP=9$, $FN=340$:
sensitivity $= 0{,}073$ và specificity $= 0{,}997$. Bắt được có $7{,}3\%$ thư rác
nhưng gần như không chặn nhầm thư thật --- với lọc thư rác đó là đánh đổi hợp lý.
Hàm tiện ích $U = TP - 50\,FP - 5\,FN$ lượng hoá đúng cái đánh đổi đó, cho
$U = 1422$.

\vd Bayes với tỉ lệ nền. Bệnh lupus có tỉ lệ mắc $2\%$, xét nghiệm có độ nhạy
$0{,}98$ và độ đặc hiệu $0{,}74$. Với $B$ là có bệnh và $+$ là dương tính:
\[
P(B \cap +) = 0{,}02 \times 0{,}98 = 0{,}0196
\qquad
P(\bar B \cap +) = 0{,}98 \times 0{,}26 = 0{,}2548
\]
\[
P(B \mid +) = \frac{0{,}0196}{0{,}0196+0{,}2548} = 0{,}0714
\]
Xét nghiệm rất nhạy nhưng dương tính chỉ cho $7\%$ khả năng mắc bệnh thật, vì bệnh
hiếm nên nhóm không bệnh đông tới mức $26\%$ dương tính giả của họ vẫn áp đảo.
\emph{Tỉ lệ nền quyết định tất cả.}

%=====================================================================
\section{Monte Carlo (chương 11)}

Thiên về mô tả thuật toán hơn là tính tay.

\subsection*{Ba trụ cột}

\textbf{Số ngẫu nhiên} --- ba loại: thực (hiện tượng vật lý), gần ngẫu nhiên (dãy
tất định phủ đều), giả ngẫu nhiên (sinh bằng thuật toán, thứ máy tính thật sự
dùng). Bốn tiêu chí của dãy tốt: chu kỳ dài, phân bố đều, độc lập thống kê, sinh
nhanh. Số giả ngẫu nhiên \emph{hoàn toàn tất định} --- cùng hạt giống thì cùng dãy,
và đó là \emph{ưu điểm} vì làm kết quả tái lập được.

\textbf{Luật số lớn} nói mô phỏng \emph{sẽ đúng}. \textbf{Định lý giới hạn trung
tâm} nói \emph{sai bao nhiêu}. Cần cả hai.

\subsection*{Các thuật toán}

\textbf{Bộ sinh đồng dư tuyến tính:} $x_i = (a\,x_{i-1}+c) \bmod m$, rồi chuẩn hoá
$x_i/m$. Chu kỳ tối đa là $m$, nên $m$ thực tế lấy $2^{32}$ trở lên.

\textbf{Box--Muller} sinh chuẩn: từ $u_1,u_2 \sim \mathcal{U}(0,1)$ đặt
$\theta = 2\pi u_1$ và $r = \sqrt{-2\ln u_2}$, trả về $r\cos\theta$ và
$r\sin\theta$ --- \emph{hai} biến $\mathcal{N}(0,1)$ độc lập.

\textbf{Biến đổi nghịch đảo:} nếu $u \sim \mathcal{U}(0,1)$ thì $X = F^{-1}(u)$ có
phân phối $F$. Rời rạc thì chọn $X = x_i$ khi $F(x_{i-1}) \le u < F(x_i)$; liên tục
thì giải $F(x)=u$. Với phân phối mũ, $x = -\ln(1-u)/\lambda$.

\textbf{Chấp nhận -- bác bỏ:} cần $q$ dễ lấy mẫu và hằng số $M$ sao cho
$p(x) \le Mq(x)$ với mọi $x$. Sinh $x \sim q$, sinh $u \sim \mathcal{U}(0, Mq(x))$,
chấp nhận nếu $u < p(x)$. Hiệu suất là $1/M$ --- chọn $q$ dở thì rất lãng phí.

\textbf{Lấy mẫu quan trọng:}
$\mathbb{E}_{p}[f] = \mathbb{E}_{q}\!\left[f\,\dfrac{p}{q}\right]$.
Lấy mẫu từ $q$ nhưng vẫn tính đúng kỳ vọng dưới $p$ nhờ trọng số $p/q$. Dùng khi
biến cố cần đo hiếm dưới $p$.

\subsection*{Tích phân Monte Carlo}

\[
\int_a^b f(x)\,dx \approx (b-a)\cdot\frac{1}{n}\sum_{i=1}^{n} f(u_i)
\qquad
SE = \frac{\sigma_f}{\sqrt n}
\]

\textbf{Hệ quả quan trọng nhất của cả chương:} sai số giảm theo $1/\sqrt n$ và
\emph{không phụ thuộc số chiều}. Cầu phương lưới cần $n^d$ điểm với $d$ chiều nên
bùng nổ tổ hợp; Monte Carlo thì không. Ở một hai chiều thì Monte Carlo \emph{thua}
hình thang, nhưng từ khoảng $d \ge 4$ thì thắng. Muốn giảm sai số một nửa phải
tăng $n$ \emph{gấp bốn}.

\vd Ước lượng $\pi$: sinh $n$ điểm ngẫu nhiên trong hình vuông, đếm $k$ điểm rơi
vào đường tròn nội tiếp, khi đó $\pi \approx 4k/n$. Số 4 đến từ chỗ hình vuông cạnh
2 có diện tích 4 còn tròn bán kính 1 có diện tích $\pi$.

\subsection*{Bài tập chương}

\textbf{Bài 1 --- thời gian dự án.} Ba công việc tuần tự với ước lượng tam giác
$(4;6;8)$, $(5;7;9)$, $(2;4;6)$ tuần. Hỏi xác suất xong trong 17 tuần. Mô phỏng:
lặp $N$ lần, mỗi lần sinh ba giá trị tam giác rồi đếm số lần tổng không quá 17.

\emph{Kiểm chứng bằng giải tích trước khi tin mô phỏng.} Phân phối tam giác có
trung bình $(a+m+b)/3$ và phương sai $(a^2+m^2+b^2-am-ab-mb)/18$. Ba công việc có
trung bình $6$, $7$, $4$ nên tổng đúng bằng \textbf{17}; phương sai mỗi cái đều là
$0{,}667$ nên tổng là $2{,}00$ với $\sigma = 1{,}41$. Cả ba đối xứng nên tổng cũng
đối xứng quanh 17, cho $P(T \le 17) \approx \mathbf{0{,}50}$. Mô phỏng chạy đúng
phải ra khoảng $0{,}50$; lệch xa là code sai.

\textbf{Bài 2 --- doanh thu thuê bao.} Phí 10 USD mỗi thuê bao mỗi tháng, 10.000
người dùng, chi phí duy trì 1.000.000 USD một năm. Giả định mỗi tháng trong 12
tháng tới giống một tháng rút ngẫu nhiên từ 60 tháng lịch sử. Điểm cần nhận ra:
bài này \emph{không giả định phân phối nào cả}, nó lấy mẫu có hoàn lại thẳng từ dữ
liệu lịch sử --- đó chính là \emph{bootstrap}.

\subsection*{Bootstrap}

Từ mẫu $n$ quan sát, rút \emph{có hoàn lại} đúng $n$ quan sát, tính lại thống kê,
lặp $B$ lần (2000--5000). Độ lệch chuẩn của $B$ giá trị đó chính là sai số chuẩn;
phân vị $2{,}5\%$ và $97{,}5\%$ cho khoảng tin cậy 95\%.

Ba câu hay bị hỏi lại: (1)~``rút hết rổ thì khác gì nhau'' --- vì rút \emph{có hoàn
lại}, nên có món trúng hai ba lần, có món không trúng lần nào;
(2)~xác suất một món không bao giờ được rút là $(1-1/n)^n \to e^{-1} \approx 0{,}368$,
nên mỗi vòng chỉ chứa khoảng \textbf{63,2\%} số món riêng biệt;
(3)~sai số chuẩn đến từ chính việc lặp lại ngẫu nhiên đó, không cần công thức giải
tích nào.

%=====================================================================
\section{Bảng tra nhanh}

\subsection*{Giá trị tới hạn}

\begin{center}
\begin{tabular}{@{}lcccc@{}}
\toprule
\textbf{Độ tin cậy} / $\alpha$ & 90\% / 0,10 & 95\% / 0,05 & 98\% / 0,02 & 99\% / 0,01 \\
\midrule
$z^{*}$ (hai phía) & 1,645 & \textbf{1,96} & 2,33 & 2,576 \\
$z_\alpha$ (một phía) & 1,282 & \textbf{1,645} & 2,054 & 2,326 \\
\bottomrule
\end{tabular}
\end{center}

$t^{*}_{0{,}025}$ theo $df$: $3 \to 3{,}182$; $5 \to 2{,}571$; $10 \to 2{,}228$;
$20 \to 2{,}086$; $25 \to 2{,}060$; $30 \to 2{,}042$; $\infty \to 1{,}96$.

$\chi^2_{0{,}05}$ theo $df$: $1 \to 3{,}841$; $2 \to 5{,}991$; $3 \to 7{,}815$;
$4 \to 9{,}488$.

Thiếu bảng tra thì ba số \textbf{1,96}, \textbf{1,645}, \textbf{2,576} gỡ được phần
lớn bài.

\subsection*{Bậc tự do dùng ở đâu}

\begin{center}
\begin{tabular}{@{}ll@{}}
\toprule
$n-1$ & phương sai mẫu; $t$ một trung bình; $\chi^2$ một phương sai \\
$n-2$ & hồi quy đơn (đã ước lượng hai tham số $b_0$ và $b_1$) \\
$n_1+n_2-2$ & $t$ hai trung bình phương sai bằng nhau \\
$(h-1)(k-1)$ & chi bình phương độc lập \\
$2$ & Jarque--Bera, luôn luôn \\
\bottomrule
\end{tabular}
\end{center}

\subsection*{Sáu câu dễ mất điểm nhất}

\begin{enumerate}
\item p-value \emph{không phải} xác suất $H_0$ đúng.
\item Không viết ``chấp nhận $H_0$'' --- viết \emph{chưa đủ bằng chứng để bác bỏ}.
\item Phương sai mẫu chia $n-1$; hồi quy dùng $df = n-2$.
\item Kiểm định một tỉ lệ dùng $p_0$ ở mẫu số; dựng khoảng tin cậy dùng $\hat p$.
\item Hồi quy trên dữ liệu quan sát thì viết \emph{đi kèm}, không viết \emph{làm cho}.
\item Điều kiện hoá collider tạo ra liên hệ giả --- thêm biến kiểm soát không phải
      lúc nào cũng an toàn hơn.
\end{enumerate}

\end{document}
